\documentclass[12pt,a4paper]{article}
\usepackage{ucebnice}

\title{Lineární funkce s absolutní hodnotou}
\author{Ondřej Škubala}
\date{\today}

\begin{document}

\maketitle
\newpage



\section{Absolutní hodnota}

\begin{quotebox}
{„Nemůžeš-li vyřešit zadaný problém, zkus nejprve vyřešit nějaký s ním související problém.“}
{George Pólya, \textit{Jak to vyřešit (1945)}}
\end{quotebox}


\begin{historicalbox}
Pojem absolutní hodnoty, který dnes působí samozřejmě a samozřejmě se používá, 
má ve skutečnosti velmi zajímavou historii.  
Jeho kořeny sahají až do doby, kdy se lidé teprve učili rozlišovat mezi „velikostí“ a „znaménkem“ čísla.

\medskip
Ve starověku se počítalo především s délkami, plochami a objemy – tedy s veličinami, které nemohou být záporné.  
Záporná čísla tehdy ještě neexistovala, a proto nebylo potřeba vyjadřovat pojem vzdálenosti zvlášť.  
Teprve v 16. a 17. století, s rozvojem algebry, obchodu a astronomie, začali matematici chápat, 
že číslo může mít i směr, a vznikla potřeba zachytit jeho „čistou velikost“ bez znaménka.  
Tento krok učinili mimo jiné \textbf{René Descartes} a \textbf{John Wallis}, kteří zavedli pojem \textit{magnituda} – 
velikost čísla nezávislá na směru.

\medskip
Současný zápis pomocí svislých čar $|x|$ se objevil až v 19. století u německého matematika 
\textbf{Karla Weierstrasse}.  
Ten hledal jednotný způsob, jak vyjádřit vzdálenost v číselné i komplexní rovině, 
a položil tím základ moderní analýzy.  
Dnes je symbol $|x|$ natolik univerzální, že jej najdete nejen v matematice, 
ale i ve fyzice či informatice – například ve funkci \texttt{abs(x)}, která počítá „velikost čísla bez směru“.

\medskip
Pokud byste si chtěli o historii matematiky přečíst víc, 
doporučuji knihu \textit{E. T. Bell – Muži, kteří dělali matematiku}, 
která poutavě popisuje příběhy velkých objevů i lidí, kteří za nimi stáli.  
Zkuste si při čtení všimnout, jak se v různých obdobích mění pohled na něco tak jednoduchého, 
jako je „vzdálenost od nuly“ –  
a přesto zůstává jednou z nejzákladnějších myšlenek celé matematiky.
\end{historicalbox}



Myšlenka, že složitější úlohu lze pochopit rozložením na jednodušší části,
je v matematice velmi užitečná.  
Podobně budeme nyní zkoumat funkce, 
které se sice na první pohled zdají jiné než běžné lineární funkce, 
ale ve skutečnosti se skládají ze dvou lineárních větví.

Než však přistoupíme k samotné definici, připomeňme si pojem \important{absolutní hodnota}, na kterém bude celá kapitola postavena.

\begin{definitionbox}
\textbf{Absolutní hodnota} reálného čísla $x$ vyjadřuje jeho vzdálenost od nuly na číselné ose.  
Označuje se symbolem $|x|$ a je definována předpisem:
\[
|x| =
\begin{cases}
x, & \text{je-li } x \ge 0,\\[4pt]
-x, & \text{je-li } x < 0.
\end{cases}
\]
Z toho plyne, že absolutní hodnota je vždy nezáporné číslo,
 tedy $|x| \ge 0$ pro všechna $x \in \mathbb{R}$.
\end{definitionbox}


Absolutní hodnota má velmi přirozený \textbf{geometrický význam}.  
Na číselné ose totiž udává \textbf{vzdálenost bodu} od počátku.  
Pokud má bod souřadnici $x$, pak délka úsečky mezi tímto bodem $X(x)$ a počátkem $O(0)$ je právě $|x|$.  

\begin{center}
\begin{tikzpicture}[scale=0.9]
    \draw[->] (-4,0)--(4,0) node[below]{$x$};
    \draw (0,0.1)--(0,-0.1) node[below]{$0$};
    \draw (2,0.1)--(2,-0.1) node[below]{$x$};
    \filldraw (2,0) circle (1.5pt);
    \draw[<->,blue!60!black,thick] (0,0.5)--(2,0.5)
        node[midway,above,black]{$|x|$};
\end{tikzpicture}
\end{center}

Tento význam lze snadno zobecnit i na dva libovolné body $A(a)$ a $B(b)$ na číselné ose.  
Jejich vzdálenost je dána rozdílem souřadnic, přičemž absolutní hodnota zaručuje,  
že výsledek bude vždy nezáporný:
\[
AB = |a - b|.
\]
Například: vzdálenost bodů $A(-3)$ a $B(5)$ je $|{-3}-5| = 8$.

Myšlenka „vzdálenosti“ pomocí absolutní hodnoty se objevuje i v dalších oblastech matematiky.  
U \textbf{komplexních čísel} $z = a + bi$ udává absolutní hodnota vzdálenost bodu $Z(a,b)$ od počátku souřadnic:
\[
|z| = \sqrt{a^2 + b^2}.
\]
Stejný princip se rozšiřuje i v \textbf{analytické geometrii},  
kde vzdálenost dvou bodů $A(x_1, y_1)$ a $B(x_2, y_2)$ vyjádříme vzorcem
\[
AB = \sqrt{(x_2 - x_1)^2 + (y_2 - y_1)^2}.
\]
Vidíme tedy, že absolutní hodnota představuje první a nejjednodušší formu pojmu \textbf{vzdálenosti} v matematice.



\begin{examplebox}
Ukažme si několik jednoduchých příkladů, jak absolutní hodnota funguje:

\[
\begin{aligned}
|5{,}1| &= 5{,}1, \\
|-0{,}8| &= 0{,}8, \\
|-7{,}2| &= 7{,}2, \\
|0| &= 0.
\end{aligned}
\]

Vidíme, že absolutní hodnota v podstatě „odstraňuje“ znaménko mínus z čísla. 
Což můžeme interpretovat jako jeho vzdálenost od nuly na číselné ose,
 bez ohledu na to, zda leží vlevo nebo vpravo od nuly.
\end{examplebox}

\subsubsection*{Pravidla pro výpočty s absolutní hodnotou}

\[
\begin{array}{cccc}
|a| \ge 0 
& \quad |a| = |-a| 
& \quad |ab| = |a| \cdot |b| 
& \quad b \ne 0 \Rightarrow \left|\frac{a}{b}\right| = \frac{|a|}{|b|}
\end{array}
\]

\smallskip

Tato pravidla nám umožňují jednoduše pracovat s absolutní hodnotou při úpravách výrazů a řešení rovnic.  
Zejména rovnost $|a| = |-a|$ ukazuje, že znaménko uvnitř absolutní hodnoty nehraje roli — vždy dostáváme stejný výsledek.

\newpage
\section{Funkce aboslutní hodnoty}
Pro každé reálné číslo $x \in \mathbb{R}$ platí podle definice, že $|x| \ge 0$.  
Je-li $x \ne 0$, pak dokonce $|x| > 0$.  
Každému reálnému číslu je tedy jednoznačně přiřazena jeho absolutní hodnota — vždy nezáporné číslo představující vzdálenost od nuly.  
Tím získáváme \important{funkci absolutní hodnoty}, kterou lze vyjádřit předpisem
$
y = |x|.
$

Jedná se tedy o funkci definovanou na celé množině reálných čísel $\mathbb{R}$, jejíž hodnota je vždy nezáporná.


Nyní si ukážeme, jak lze funkci absolutní hodnoty chápat jako spojení dvou jednoduchých lineárních funkcí.

\begin{examplebox}
Abychom lépe pochopili tvar grafu funkce absolutní hodnoty, začněme dvěma základními lineárními funkcemi.

\begin{center}
\begin{minipage}{0.45\textwidth}
\centering
\begin{axisgraph}[
    xmin=-4, xmax=4,
    ymin=-3, ymax=3,
    axis lines=middle,
    xlabel={$x$}, ylabel={$y$},
    title={$y = x$},
    width=\linewidth,
    height=5.5cm,
    samples=100
]
    \addplot[thick, pastelblue!80!black, domain=-3.5:3.5] {x};
\end{axisgraph}
\end{minipage}
\hfill
\begin{minipage}{0.45\textwidth}
\centering
\begin{axisgraph}[
    xmin=-4, xmax=4,
    ymin=-3, ymax=3,
    axis lines=middle,
    xlabel={$x$}, ylabel={$y$},
    title={$y = -x$},
    width=\linewidth,
    height=5.5cm,
    samples=100
]
    \addplot[thick, pastelblue!80!black, domain=-3.5:3.5] {-x};
\end{axisgraph}
\end{minipage}
\end{center}

Funkce $y = x$ má rostoucí charakter a prochází počátkem souřadnic.  
Naopak funkce $y = -x$ je klesající a je k funkci $y = x$ souměrná podle osy $x$.  
Pokud z těchto dvou funkcí ponecháme pouze části nad osou $x$ (neboli jejich kladné hodnoty) a tyto části spojíme, vznikne nový graf, který nazýváme \important{funkce absolutní hodnoty}.

\begin{center}
\begin{axisgraph}[
    xmin=-4, xmax=4,
    ymin=-1, ymax=4,
    axis lines=middle,
    xlabel={$x$}, ylabel={$y$},
    title={Funkce $y = |x|$},
    width=0.8\linewidth,
    height=6cm,
    samples=100
]
    % levá a pravá větev absolutní hodnoty
    \addplot[thick, pastelblue!80!black, domain=0:3.5] {x};
    \addplot[thick, pastelblue!80!black, domain=-3.5:0] {-x};
    % vrchol
    \addplot[only marks, mark=*, pastelblue!80!black] coordinates {(0,0)};
\end{axisgraph}
\end{center}

Vidíme, že graf funkce $y = |x|$ vznikl „překlopením“ části grafu $y = x$, která ležela pod osou $x$, směrem nahoru.  
Vrchol této funkce leží v počátku souřadnic, a právě zde se graf „láme“ – mění směr, ale zůstává spojitý.
\end{examplebox}


\subsection{Obecný tvar funkce s absolutní hodnotou}

Nyní zobecníme poznatky o základní funkci $y = |x|$ na funkce, které obsahují výraz $a x + b$ uvnitř absolutní hodnoty a mohou být navíc posunuty po ose $y$.

\begin{definitionbox}
\textbf{Lineární funkcí s absolutní hodnotou} nazýváme každou funkci tvaru
\[
f(x) = |a x + b| + c,
\]
kde $a, b, c \in \mathbb{R}$ a $a \ne 0$.
\end{definitionbox}

Taková funkce se skládá ze dvou lineárních částí – podle toho, zda je výraz $a x + b$ kladný, nebo záporný.  
Použijeme-li definici absolutní hodnoty, můžeme zápis přepsat do tvaru:
\[
|a x + b| =
\begin{cases}
a x + b, & \text{je-li } a x + b \ge 0,\\[4pt]
-(a x + b), & \text{je-li } a x + b < 0.
\end{cases}
\]

Po dosazení do původního předpisu dostaneme:
\[
f(x) =
\begin{cases}
a x + b + c, & \text{je-li } a x + b \ge 0,\\[6pt]
- a x - b + c, & \text{je-li } a x + b < 0.
\end{cases}
\]

Získali jsme tedy dvě přímky, které se liší pouze znaménkem u členů s $a$ a $b$.  
Otázkou zůstává, kde přesně se tyto dvě větve spojí.

\subsubsection*{Odvození vzorce pro vrchol absolutní hodnoty}

Víme, že k \textit{„překlápění“} dochází tehdy, když se mění znaménko výrazu uvnitř absolutní hodnoty – tedy výrazu $a x + b$.  
Pokud je tento výraz kladný, absolutní hodnota ho nijak nezmění; pokud je záporný, změní se na opačné znaménko.  
Přechod mezi těmito dvěma situacemi nastává právě ve chvíli, kdy výraz $a x + b$ přechází ze záporné hodnoty do kladné, tedy když se rovná nule.

\[
a x + b = 0
\]

Tento bod je hranicí mezi levou a pravou částí grafu a určuje místo, kde se funkce \textit{„láme“}.  
Z uvedené rovnice snadno vyjádříme souřadnici $x$ vrcholu:

\begin{align*}
a x + b &= 0 \\[4pt]
a x &= -b \\[4pt]
x_v &= -\frac{b}{a}
\end{align*}

Pro tuto hodnotu proměnné je výraz uvnitř absolutní hodnoty nulový, tedy $|a x_v + b| = 0$, a hodnota funkce je proto rovna konstantě $c$:
\[
f(x_v) = c.
\]

Souřadnice vrcholu funkce jsou tedy
\[
V\left[-\frac{b}{a};\, c\right].
\]

\begin{definitionbox}
\textbf{Vrchol lineární funkce s absolutní hodnotou} je bod, v němž se graf funkce láme – tedy místo, kde se výraz $a x + b$ uvnitř absolutní hodnoty rovná nule.  
Jeho souřadnice jsou
\[
V\left[-\frac{b}{a};\, c\right],
\]
kde $a, b, c \in \mathbb{R}$ a $a \ne 0$.
\end{definitionbox}



\subsection{Graf lineární funkce s absolutní hodnotou}

Dříve než si vysvětlíme, jak jednotlivé parametry $a$, $b$ a $c$ ovlivňují tvar a polohu grafu, podívejme se na několik ukázek různých funkcí s absolutní hodnotou:

\begin{center}
\begin{minipage}{0.32\textwidth}
\centering
\begin{axisgraph}[
    xmin=-2, xmax=6,
    ymin=-1, ymax=5,
    axis lines=middle,
    xlabel={$x$}, ylabel={$y$},
    title={$y = |x - 2|$},
    width=\linewidth,
    height=5.5cm,
    samples=100
]
    \addplot[thick, pastelblue!80!black, domain=-2:2] {-x + 2};
    \addplot[thick, pastelblue!80!black, domain=2:6] {x - 2};
\end{axisgraph}
\end{minipage}
\hfill
\begin{minipage}{0.32\textwidth}
\centering
\begin{axisgraph}[
    xmin=-2, xmax=6,
    ymin=-1, ymax=5,
    axis lines=middle,
    xlabel={$x$}, ylabel={$y$},
    title={$y = |x - 3| + 1$},
    width=\linewidth,
    height=5.5cm,
    samples=100
]
    \addplot[thick, pastelblue!80!black, domain=-2:3] {-x + 4};
    \addplot[thick, pastelblue!80!black, domain=3:6] {x - 2};
\end{axisgraph}
\end{minipage}
\hfill
\begin{minipage}{0.32\textwidth}
\centering
\begin{axisgraph}[
    xmin=-4, xmax=3,
    ymin=-4, ymax=4,
    axis lines=middle,
    xlabel={$x$}, ylabel={$y$},
    title={$y = 2|x + 1| - 3$},
    width=\linewidth,
    height=5.5cm,
    samples=100
]
    \addplot[thick, pastelblue!80!black, domain=-4:-1] {-2*(x + 1) - 3};
    \addplot[thick, pastelblue!80!black, domain=-1:3] {2*(x + 1) - 3};
\end{axisgraph}
\end{minipage}
\end{center}

\vspace{1em}

Na první pohled vidíme, že všechny tři grafy mají stejný tvar písmene „V“, ale liší se svou polohou, výškou i strmostí.  
V následující části si krok za krokem vysvětlíme, jak jednotlivé parametry $a$, $b$ a $c$ ovlivňují tyto změny.


\subsubsection*{Posun grafu podle parametru $b$}

Parametr $b$ určuje, \textit{kde na ose $x$} se bude nacházet vrchol funkce.  
Základní tvar funkce je
\[
y = |x - b|.
\]
Hodnota $b$ ovlivňuje posun grafu vodorovným směrem.  
Podle znaménka výrazu uvnitř absolutní hodnoty se graf posouvá následovně:

\begin{itemize}
    \item Pokud od $x$ \textbf{odečítáme} číslo $b$, graf se posune o $b$ jednotek doprava.  
    \item Pokud k $x$ \textbf{přičítáme} číslo $b$, graf se posune o $b$ jednotek doleva. 
\end{itemize}

\begin{examplebox}
Porovnejme dvě funkce:
\[
y_1 = |x|, \qquad y_2 = |x - 3|.
\]
Funkce $y_1$ má vrchol v bodě $[0;0]$, zatímco funkce $y_2$ v bodě $[3;0]$.  
Vrchol se tedy posunul o 3 jednotky \textbf{doprava}.

\begin{center}
\begin{axisgraph}[
    xmin=-2, xmax=6,
    ymin=-1, ymax=4,
    axis lines=middle,
    xlabel={$x$}, ylabel={$y$},
    title={Posun grafu podle $b$},
    width=0.7\linewidth,
    height=6cm,
    samples=100
]
    % původní funkce y=|x|
    \addplot[thick, gray!85!black, domain=-2:0] {-x};
    \addplot[thick, gray!85!black, domain=0:6] {x};
    % posunutá funkce y=|x-3|
    \addplot[thick, pastelorange!80!black, domain=-2:3] {-x + 3};
    \addplot[thick, pastelorange!80!black, domain=3:6] {x - 3};
    % vrcholy
    \addplot[only marks, mark=*, gray!85!black] coordinates {(0,0)};
    \addplot[only marks, mark=*, pastelorange!80!black] coordinates {(3,0)};
    % popisky
    \node[below right, gray!70] at (axis cs:0,0) {$[0;0]$};
    \node[above, pastelorange!80!black] at (axis cs:3,0.5) {$[3;0]$};
\end{axisgraph}
\end{center}
\end{examplebox}

\begin{notebox}
Vrchol funkce $y = |x - b|$ má souřadnice $V[b;0]$.  
Pokud $b > 0$, graf se posune doprava; pokud $b < 0$, posune se doleva.
\end{notebox}



\subsubsection*{Vliv parametru $a$ na tvar grafu}

Parametr $a$ ovlivňuje, jak je graf \textit{otevřený} (roztažený či zúžený) a zároveň, 
zda směřuje \textit{nahoru} nebo \textit{dolů}.  
Zapisujeme funkci ve tvaru
\[
y = |a x|.
\]

\begin{itemize}
  \item Pokud je $|a| > 1$, graf je \important{strmější} — větve jsou blíže k ose $y$.  
  \item Pokud je $0 < |a| < 1$, graf je \important{plošší} — větve jsou „roztažené“ do stran.  
  \item Pokud je $a < 0$, graf se \important{převrátí směrem dolů}.
\end{itemize}

\begin{examplebox}
Porovnejme tři funkce:
\[
y_1 = |x|, \quad y_2 = |2x|, \quad y_3 = |{\tfrac{1}{2}}x|.
\]
\begin{center}
\begin{minipage}{0.32\textwidth}
\centering
\begin{axisgraph}[
    xmin=-4, xmax=4,
    ymin=-1, ymax=4,
    axis lines=middle,
    xlabel={$x$}, ylabel={$y$},
    title={$y = |x|$},
    width=\linewidth,
    height=5.5cm,
    samples=100
]
    \addplot[thick, gray!80!black, domain=-4:0] {-x};
    \addplot[thick, gray!80!black, domain=0:4] {x};
\end{axisgraph}
\end{minipage}
\hfill
\begin{minipage}{0.32\textwidth}
\centering
\begin{axisgraph}[
    xmin=-4, xmax=4,
    ymin=-1, ymax=4,
    axis lines=middle,
    xlabel={$x$}, ylabel={$y$},
    title={$y = |2x|$},
    width=\linewidth,
    height=5.5cm,
    samples=100
]
    \addplot[thick, pastelblue!80!black, domain=-4:0] {-2*x};
    \addplot[thick, pastelblue!80!black, domain=0:4] {2*x};
\end{axisgraph}
\end{minipage}
\hfill
\begin{minipage}{0.32\textwidth}
\centering
\begin{axisgraph}[
    xmin=-4, xmax=4,
    ymin=-1, ymax=4,
    axis lines=middle,
    xlabel={$x$}, ylabel={$y$},
    title={$y = \tfrac{1}{2}|x|$},
    width=\linewidth,
    height=5.5cm,
    samples=100
]
    \addplot[thick, pastelorange!80!black, domain=-4:0] {-(x/2)};
    \addplot[thick, pastelorange!80!black, domain=0:4] {x/2};
\end{axisgraph}
\end{minipage}
\end{center}


Graf $y = |2x|$ je dvakrát strmější, zatímco $y = \tfrac{1}{2}|x|$ je dvakrát plošší než základní funkce $y = |x|$.
\end{examplebox}

\begin{notebox}
\textbf{Shrnutí vlivu parametru $a$:}
\begin{itemize}
    \item $a > 1$ – zúžení (graf strmější),
    \item $0 < a < 1$ – roztažení (graf plošší),
    \item $a < 0$ – převrácení směrem dolů,
    \item $|a|$ – určuje míru „otevřenosti“ tvaru.
\end{itemize}
\end{notebox}

\subsubsection*{Rozdíl mezi $y = |a x|$ a $y = a |x|$}

Je důležité rozlišovat, zda se koeficient $a$ nachází \textit{uvnitř} nebo \textit{vně} absolutní hodnoty.

\[
y = |a x| \qquad \text{a} \qquad y = a|x|
\]

\begin{itemize}
  \item U funkce $y = |a x|$ se absolutní hodnota vztahuje i na $a$, takže
  \[
  |a x| = |a| \cdot |x|,
  \]
  a graf je vždy otevřen směrem nahoru – změna znaménka $a$ zde nic nemění.
  
  \item U funkce $y = a|x|$ se znaménko $a$ uplatní \textit{až po} výpočtu absolutní hodnoty.  
  Proto když $a < 0$, celý graf se převrátí dolů.
\end{itemize}

\begin{center}
\begin{axisgraph}[
    xmin=-4, xmax=4,
    ymin=-4, ymax=4,
    axis lines=middle,
    xlabel={$x$}, ylabel={$y$},
    title={Rozdíl mezi $|a x|$ a $a|x|$},
    width=0.8\linewidth,
    height=6cm,
    samples=100
]
    % y=|2x|
    \addplot[thick, pastelblue!80!black, domain=-4:0] {-2*x};
    \addplot[thick, pastelblue!80!black, domain=0:4] {2*x};
    % y=-2|x|
    \addplot[thick, pastelorange!80!black, domain=-4:0] {2*x};
    \addplot[thick, pastelorange!80!black, domain=0:4] {-2*x};
    % popisky
    \node[above right, pastelblue!80!black] at (axis cs:1,3) {$y=|2x|$};
    \node[below right, pastelorange!80!black] at (axis cs:1,-3) {$y=-2|x|$};
\end{axisgraph}
\end{center}

\begin{notebox}
Rozdíl mezi $y = |a x|$ a $y = a|x|$ spočívá v tom, že v prvním případě se 
znaménko $a$ \textit{ztrácí} uvnitř absolutní hodnoty,  
zatímco v druhém případě zůstává zachováno a ovlivňuje orientaci celého grafu.
\end{notebox}




\subsubsection*{Posun grafu podle parametru $c$}

Parametr $c$ ovlivňuje, \textit{kde na ose $y$} se nachází vrchol funkce.  
V základním tvaru
\[
y = |x| + c
\]
představuje číslo $c$ \important{svislý posun} celého grafu – nemění se jeho tvar ani sklon větví, pouze poloha.

\begin{itemize}
    \item Pokud je $c > 0$, graf se posune \important{nahoru} o $c$ jednotek.  
    \item Pokud je $c < 0$, graf se posune \important{dolů} o $|c|$ jednotek.  
\end{itemize}

\begin{examplebox}
Porovnejme tři funkce:
\[
y_1 = |x|, \quad y_2 = |x| + 2, \quad y_3 = |x| - 3.
\]

\begin{center}
\begin{minipage}{0.32\textwidth}
\centering
\begin{axisgraph}[
    xmin=-4, xmax=4,
    ymin=-4, ymax=4,
    axis lines=middle,
    xlabel={$x$}, ylabel={$y$},
    title={$y = |x|$},
    width=1.25\linewidth,
    height=7cm,
    samples=100
]
    \addplot[thick, gray!80!black, domain=-4:0] {-x};
    \addplot[thick, gray!80!black, domain=0:4] {x};
\end{axisgraph}
\end{minipage}
\hfill
\begin{minipage}{0.32\textwidth}
\centering
\begin{axisgraph}[
    xmin=-4, xmax=4,
    ymin=-4, ymax=4,
    axis lines=middle,
    xlabel={$x$}, ylabel={$y$},
    title={$y = |x| + 2$},
    width=1.25\linewidth,
    height=7cm,
    samples=100
]
    \addplot[thick, pastelblue!80!black, domain=-4:0] {-x + 2};
    \addplot[thick, pastelblue!80!black, domain=0:4] {x + 2};
\end{axisgraph}
\end{minipage}
\hfill
\begin{minipage}{0.32\textwidth}
\centering
\begin{axisgraph}[
    xmin=-4, xmax=4,
    ymin=-4, ymax=4,
    axis lines=middle,
    xlabel={$x$}, ylabel={$y$},
    title={$y = |x| - 3$},
    width=1.25\linewidth,
    height=7cm,
    samples=100
]
    \addplot[thick, pastelorange!80!black, domain=-4:0] {-x - 3};
    \addplot[thick, pastelorange!80!black, domain=0:4] {x - 3};
\end{axisgraph}
\end{minipage}
\end{center}

Graf $y = |x| + 2$ je posunut o 2 jednotky nahoru, zatímco graf $y = |x| - 3$ o 3 jednotky dolů.
\end{examplebox}

\begin{notebox}
    \textbf{Parametr $c$} způsobuje \textit{svislý posun} grafu:
    \begin{itemize}
        \item $c > 0$ – posun nahoru,  
        \item $c < 0$ – posun dolů.
    \end{itemize}
    Vrchol funkce má souřadnice $V[0; c]$.
\end{notebox}


\newpage
\subsubsection*{Kombinace všech parametrů}
Nyní už víme, že parametr $a$ ovlivňuje \textit{strmost a orientaci} grafu,
parametr $b$ \textit{posune po ose $x$} a parametr $c$ \textit{posune po ose $y$}.

Abychom viděli, jak všechny tři působí dohromady, podíváme se na několik příkladů
ve kterých se kombinují všechny předchozí poznatky.
Na následujících dvou funkcích budeme sledovat, jak se mění tvar i poloha grafu.

\begin{examplebox}
\textbf{Příklad:} \(y = |x - 2| - 3\)

Postupně si ukážeme, jak tento graf vzniká z funkce \(y = |x|\).

\medskip
\textbf{1. krok – základní tvar:}
\[
y = |x|
\]
Graf má vrchol v počátku $[0;0]$ a tvoří základní tvar písmene „V“.

\begin{center}
\begin{axisgraph}[
    xmin=-2, xmax=6,
    ymin=-5, ymax=3,
    axis lines=middle,
    axis line style={thick},
    xlabel={$x$}, ylabel={$y$},
    title={Základní tvar $y = |x|$},
    width=0.75\linewidth,
    height=6.5cm,
    samples=100,
    xtick={0,2,4,6},
    ytick={-5,-3,0,3},
    grid=none
]
    \addplot[thick, gray!75!black, domain=-2:0] {-x};
    \addplot[thick, gray!75!black, domain=0:6] {x};
    \addplot[only marks, mark=*, gray!75!black] coordinates {(0,0)};
    \node[below right, gray!70] at (axis cs:0,0) {$[0;0]$};
\end{axisgraph}
\end{center}

\medskip
\textbf{2. krok – posun o $b=2$ doprava:}
\[
y = |x - 2|
\]
Protože od $x$ odečítáme číslo 2, celý graf posuneme o 2 jednotky doprava.  
Vrchol se posune do bodu $[2;0]$.

\begin{center}
\begin{axisgraph}[
    xmin=-2, xmax=6,
    ymin=-5, ymax=3,
    axis lines=middle,
    axis line style={thick},
    xlabel={$x$}, ylabel={$y$},
    title={Posun o $b = 2$ doprava},
    width=0.75\linewidth,
    height=6.5cm,
    samples=100,
    xtick={0,2,4,6},
    ytick={-5,-3,0,3},
    grid=none
]
    \addplot[thick, gray!70!black, domain=-2:0] {-x};
    \addplot[thick, gray!70!black, domain=0:6] {x};
    \addplot[thick, pastelblue!80!black, domain=-2:2] {-x + 2};
    \addplot[thick, pastelblue!80!black, domain=2:6] {x - 2};
    \addplot[only marks, mark=*, pastelblue!80!black] coordinates {(2,0)};
    \node[below right, pastelblue!80!black] at (axis cs:2,0) {$[2;0]$};
\end{axisgraph}
\end{center}

\medskip
\textbf{3. krok – posun o $c=-3$ dolů:}
\[
y = |x - 2| - 3
\]
Hodnota $c = -3$ posune celý graf o tři jednotky dolů.  
Vrchol se tak přesune do bodu $[2;-3]$.

\begin{center}
\begin{axisgraph}[
    xmin=-2, xmax=6,
    ymin=-5, ymax=3,
    axis lines=middle,
    axis line style={thick},
    xlabel={$x$}, ylabel={$y$},
    title={Konečný tvar $y = |x - 2| - 3$},
    width=0.75\linewidth,
    height=6.5cm,
    samples=100,
    xtick={0,2,4,6},
    ytick={-5,-3,0,3},
    grid=none
]
    \addplot[thick, gray!70!black, domain=-2:2] {-x + 2};
    \addplot[thick, gray!70!black, domain=2:6] {x - 2};
    \addplot[thick, pastelorange!80!black, domain=-2:2] {-x - 1};
    \addplot[thick, pastelorange!80!black, domain=2:6] {x - 5};
    \addplot[only marks, mark=*, pastelorange!80!black] coordinates {(2,-3)};
    \node[below right, pastelorange!80!black] at (axis cs:2,-3) {$[2;-3]$};
\end{axisgraph}
\end{center}

Vrchol se posunul o 2 jednotky doprava a o 3 jednotky dolů.  
Směr i sklon větví zůstaly stejné jako u základní funkce $y = |x|$.
\end{examplebox}



Vidíme tedy, že s parametry $b$ a $c$ můžeme celý graf posouvat po rovině 
nahoru, dolů, doprava či doleva.  
Ale co se stane, pokud k tomu přidáme i násobek $a$, který mění směr a sklon grafu?  
Podívejme se na následující příklad.


\begin{examplebox}
\textbf{Příklad:} \(y = -2\,|x + 2{,}5| + 1{,}5\)

Ukážeme si krok za krokem, jak z funkce \(y = |x|\) postupně vznikne výsledný graf.

\medskip
\textbf{1. krok – základní tvar:}
\[
y = |x|
\]
Základní graf má vrchol v počátku $[0;0]$ a je otevřen směrem nahoru.

\begin{center}
\begin{axisgraph}[
    xmin=-6, xmax=2,
    ymin=-5, ymax=4,
    axis lines=middle,
    axis line style={thick},
    xlabel={$x$}, ylabel={$y$},
    title={Základní tvar $y = |x|$},
    width=0.75\linewidth,
    height=6.5cm,
    samples=100,
    xtick={-6,-4,-2,0,2},
    ytick={-5,-3,0,3},
    grid=none
]
    \addplot[thick, gray!75!black, domain=-6:0] {-x};
    \addplot[thick, gray!75!black, domain=0:2] {x};
    \addplot[only marks, mark=*, gray!75!black] coordinates {(0,0)};
    \node[below right, gray!70] at (axis cs:0,0) {$[0;0]$};
\end{axisgraph}
\end{center}

\medskip
\textbf{2. krok – posun o $b=-2{,}5$ doleva:}
\[
y = |x + 2{,}5|
\]
Protože k $x$ přičítáme číslo 2,5, posuneme celý graf o 2,5 jednotky \important{doleva}.  
Vrchol se posune do bodu $[-2{,}5;0]$.

\begin{center}
\begin{axisgraph}[
    xmin=-6, xmax=2,
    ymin=-5, ymax=4,
    axis lines=middle,
    axis line style={thick},
    xlabel={$x$}, ylabel={$y$},
    title={Posun o $b = -2{,}5$ doleva},
    width=0.75\linewidth,
    height=6.5cm,
    samples=100,
    xtick={-6,-4,-2,0,2},
    ytick={-5,-3,0,3},
    grid=none
]
    % původní
    \addplot[thick, gray!70!black, domain=-6:0] {-x};
    \addplot[thick, gray!70!black, domain=0:2] {x};
    % posunutý
    \addplot[thick, pastelblue!80!black, domain=-6:-2.5] {-(x + 2.5)};
    \addplot[thick, pastelblue!80!black, domain=-2.5:2] {(x + 2.5)};
    \addplot[only marks, mark=*, pastelblue!80!black] coordinates {(-2.5,0)};
    \node[below right, pastelblue!80!black] at (axis cs:-2.5,0) {$[-2{,}5;0]$};
\end{axisgraph}
\end{center}

\medskip
\textbf{3. krok – násobek $a=-2$:}
\[
y = -2\,|x + 2{,}5|
\]
Znaménko „$-$“ graf převrátí \important{směrem dolů} a současně hodnota $|a|=2$ graf \important{zúží} – je strmější.

Vrchol zůstává v bodu $[-2{,}5;0]$.

\begin{center}
\begin{axisgraph}[
    xmin=-6, xmax=2,
    ymin=-5, ymax=4,
    axis lines=middle,
    axis line style={thick},
    xlabel={$x$}, ylabel={$y$},
    title={Převrácení a zúžení ($a = -2$)},
    width=0.75\linewidth,
    height=6.5cm,
    samples=100,
    xtick={-6,-4,-2,0,2},
    ytick={-5,-3,0,3},
    grid=none
]
    % předchozí modrý
    \addplot[thick, gray!70!black, domain=-6:-2.5] {-(x + 2.5)};
    \addplot[thick, gray!70!black, domain=-2.5:2] {(x + 2.5)};
    % nový
    \addplot[thick, pastelorange!80!black, domain=-6:-2.5] {2*x + 5};
    \addplot[thick, pastelorange!80!black, domain=-2.5:2] {-2*x - 5};
    \addplot[only marks, mark=*, pastelorange!80!black] coordinates {(-2.5,0)};
    \node[above right, pastelorange!80!black] at (axis cs:-2.5,0) {$[-2{,}5;0]$};
\end{axisgraph}
\end{center}

\medskip
\textbf{4. krok – posun o $c = +1{,}5$ nahoru:}
\[
y = -2\,|x + 2{,}5| + 1{,}5
\]
Hodnota $c$ posouvá celý graf o 1,5 jednotky nahoru.  
Vrchol se přesune do bodu $[-2{,}5;1{,}5]$.

\begin{center}
\begin{axisgraph}[
    xmin=-6, xmax=2,
    ymin=-5, ymax=4,
    axis lines=middle,
    axis line style={thick},
    xlabel={$x$}, ylabel={$y$},
    title={Konečný tvar $y = -2|x + 2{,}5| + 1{,}5$},
    width=0.75\linewidth,
    height=6.5cm,
    samples=100,
    xtick={-6,-4,-2,0,2},
    ytick={-5,-3,0,1.5,3},
    grid=none
]
    % předchozí šedý
    \addplot[thick, gray!70!black, domain=-6:-2.5] {2*x + 5};
    \addplot[thick, gray!70!black, domain=-2.5:2] {-2*x - 5};
    % finální oranžový
    \addplot[thick, pastelorange!80!black, domain=-6:-2.5] {2*x + 6.5};
    \addplot[thick, pastelorange!80!black, domain=-2.5:2] {-2*x - 3.5};
    \addplot[only marks, mark=*, pastelorange!80!black] coordinates {(-2.5,1.5)};
    \node[above right, pastelorange!80!black] at (axis cs:-2.5,1.5) {$[-2{,}5;\,1{,}5]$};
\end{axisgraph}
\end{center}

Vrchol se tedy nachází v bodu $[-2{,}5;\,1{,}5]$,  
graf je otevřen směrem dolů a je dvakrát strmější než základní tvar $y = |x|$.
\end{examplebox}



\subsection{Funkce s více absolutními hodnotami}

Dosud jsme pracovali s funkcemi, které obsahovaly pouze jednu absolutní hodnotu.  
Nyní se podíváme na případy, kdy se v předpisu vyskytuje více absolutních hodnot.  
V takových situacích je nutné určit \important{nulové body výrazů uvnitř absolutních hodnot}  
a podle nich rozdělit osu $x$ na intervaly.  
V každém z těchto intervalů pak odstraníme absolutní hodnoty podle jejich znamének a funkci určíme zvlášť.

\begin{examplebox}
\textbf{Příklad:} \(y = |x - 2| + |x|\)

\medskip
\textbf{1. krok – určení nulových bodů:}

Každý výraz uvnitř absolutní hodnoty může měnit znaménko v bodě, kde se rovná nule:
\[
x - 2 = 0 \Rightarrow x = 2, 
\qquad x = 0.
\]
Rozdělíme tedy osu $x$ na tři intervaly:
\[
(-\infty; 0), \quad (0; 2), \quad (2; \infty).
\]

\medskip
\textbf{2. krok – odstranění absolutních hodnot v jednotlivých intervalech:}

Pomocí nulových bodů $x=0$ a $x=2$ rozdělíme reálnou osu na tři intervaly a sledujeme,  
jaké hodnoty mají výrazy uvnitř absolutních hodnot v každém z nich.

\arrayrulecolor{tableborder}
\renewcommand{\arraystretch}{1.4}
\begin{center}
\begin{tabular}{|C{2.5cm}|C{2.5cm}|C{2.5cm}|C{2.5cm}|}
\hline
\rowcolor{tableheaderdark!70} 
\textbf{Výraz} & \textbf{$(-\infty; 0)$} & \textbf{$(0; 2)$} & \textbf{$(2; \infty)$} \\ \hline
$x - 2$ & záporný & záporný & kladný \\ \hline
\(|x - 2|\) & $-(x - 2)$ & $-(x - 2)$ & $x - 2$ \\ \hline
$x$ & záporný & kladný & kladný \\ \hline
\(|x|\) & $-x$ & $x$ & $x$ \\ \hline
\textbf{$y = |x - 2| + |x|$} & $-(x - 2) - x = -2x + 2$ & $-(x - 2) + x = 2$ & $x - 2+ x = 2x - 2$ \\ \hline
\end{tabular}
\end{center}
\arrayrulecolor{black}
\renewcommand{\arraystretch}{1.0}

\medskip
Z tabulky vyčteme, že výsledný předpis má tři části:
\[
y =
\begin{cases}
-2x + 2, & x < 0,\\[4pt]
2, & 0 \le x < 2,\\[4pt]
2x - 2, & x \ge 2.
\end{cases}
\]


\medskip
\textbf{3. krok – výsledný předpis:}
\[
y =
\begin{cases}
-2x + 2, & x < 0,\\[4pt]
2, & 0 \le x < 2,\\[4pt]
2x - 2, & x \ge 2.
\end{cases}
\]

\begin{center}
\begin{axisgraph}[
    xmin=-4, xmax=4,
    ymin=-2, ymax=5,
    axis lines=middle,
    axis line style={thick},
    xlabel={$x$}, ylabel={$y$},
    title={$y = |x - 2| + |x|$},
    width=0.8\linewidth,
    height=7cm,
    samples=100,
    xtick={-4,-2,0,2,4},
    ytick={0,2,4},
    grid=none
]
    % x < 0
    \addplot[thick, pastelorange!80!black, domain=-4:0] {-2*x + 2};
    % 0 <= x < 2
    \addplot[thick, pastelblue!80!black, domain=0:2] {2};
    % x >= 2
    \addplot[thick, pastelgreen!70!black, domain=2:4] {2*x - 2};
    % body zlomu
    \addplot[only marks, mark=*, gray!70!black] coordinates {(0,2) (2,2)};
    \node[below left, gray!60] at (axis cs:0,2) {$[0;2]$};
    \node[below right, gray!60] at (axis cs:2,2) {$[2;2]$};
\end{axisgraph}
\end{center}

\textit{Graf se skládá ze tří lineárních částí, které se lámou v bodech $x=0$ a $x=2$.}
\end{examplebox}

\begin{notebox}
Při práci s více absolutními hodnotami postupujeme vždy stejně:
\begin{enumerate}[label=\alph*)]
  \item Určíme nulové body výrazů uvnitř absolutních hodnot.  
  \item Rozdělíme reálnou osu podle těchto bodů na intervaly.  
  \item V každém intervalu odstraníme absolutní hodnoty podle znaménka výrazů.  
  \item Sestavíme výsledný předpis a graf po částech.
\end{enumerate}
\end{notebox}





% -----------------------------------------------------------------
% ---------------------- P Ř Í K L A D Y --------------------------
% -----------------------------------------------------------------

\exercisepart
\setcounter{section}{1} % případně uprav dle předchozího číslování


% =========================
% ÚLOHA 1
% =========================
\begin{exercise}[Kdy platí $|a|=-a$]
    Pro která z čísel $a=\{3;\,-1{,}5;\,-13{,}4;\,2{,}7\}$ platí rovnost $|a|=-a$?
\end{exercise}


% =========================
% ÚLOHA 2
% =========================
\begin{exercise}[Základní výpočty absolutní hodnoty]
    Vypočítejte hodnoty následujících výrazů:
    \begin{tasks}(4)
        \task $|5|$
        \task $|-7|$
        \task $|0|$
        \task $|-3{,}6|$
        \task $|+12|$
        \task $|-10{,}2|$
        \task $|-0{,}1|$
        \task $|-100|$
    \end{tasks}
\end{exercise}


% =========================
% ÚLOHA 3
% =========================
\begin{exercise}[Absolutní hodnota ve výrazech]
    Vypočítejte:
    \begin{tasks}(2)
        \task $|-\sqrt{2}| - |\sqrt{2}|$
        \task $|8-11| - |3-9|$
        \task $|-3{,}5 - (-0{,}8)|$
        \task $14 - |4{,}6 - 5{,}3|$
        \task $|-2 - |5-7||$
        \task $|-|2-5| + |1| - 7|$
        \task $||-8| - |4-10||$
        \task $||2| - |3-5||$
    \end{tasks}
\end{exercise}


% =========================
% ÚLOHA 4
% =========================
\begin{exercise}[Definiční obory funkcí s absolutní hodnotou]
    Vyjádřete pomocí intervalů definiční obor těchto funkcí:
    \begin{tasks}(3)
        \task $y=\sqrt{|x|-x}$
        \task $y=\sqrt{|x|+x}$
        \task $y=\dfrac{1}{\sqrt{|x|+x}}$
    \end{tasks}
\end{exercise}


% =========================
% ÚLOHA 5
% =========================
\begin{exercise}[Základní grafy funkcí s absolutní hodnotou]
    Načrtněte grafy následujících funkcí:
    \begin{tasks}(3)
        \task $y=|x|$
        \task $y=|x-1|$
        \task $y=|x+1{,}5|$
        \task $y=|2+0{,}5x|$
        \task $y=|5-2x|$
        \task $y=|2x|$
        \task $y=|x|+4$
        \task $y=1-|0{,}5x|$
        \task $y=x+|x|$
        \task $y=x-|x|$
        \task $y=||x|-4|$
        \task $y=\dfrac{|x|}{x}$
    \end{tasks}
\end{exercise}


% =========================
% ÚLOHA 6
% =========================
\begin{exercise}[Složené funkce s absolutní hodnotou]
    Narýsujte grafy těchto funkcí:
    \begin{tasks}(2)
        \task $f_1: \; y=|x|$
        \task $f_2: \; y=|x|+2$
        \task $f_3: \; y=|x+2|$
        \task $f_4: \; y=|x+2|-3$
        \task $g_1: \; y=2|x|$
        \task $g_2: \; y=2|x-3|$
        \task $g_3: \; y=2|x-3|+1$
        \task $h_1: \; y=|x+2| + 0{,}5|x-1| - x$
        \task $h_2: \; y=|x+1| - |3-x| + 2$
        \task $t_1: \; y=|x-1| - 2$
        \task $t_2: \; y=||x-1|-2| - 3$
        \task $t_3: \; y=|||x-1|-2|-3|$
        \task $t_4: \; y=||||x|-1|-1|-1|$
        \task $r_1: \; y=|x-3| - |5-2x| + 3|1-x|$
    \end{tasks}
\end{exercise}


% =========================
% ÚLOHA 7
% =========================
\begin{exercise}[Aplikační úlohy s absolutní hodnotou]
    Absolutní hodnota vyjadřuje v mnoha reálných situacích \textbf{vzdálenost}, \textbf{změnu} či \textbf{odchylku}.  
    Vyřešte následující úlohy a výsledek vždy zapište pomocí absolutní hodnoty.

    \begin{tasks}(1)
        \task \textbf{Teplota:}  
        V poledne bylo $12\,^\circ\mathrm{C}$ a večer $-3\,^\circ\mathrm{C}$.  
        Jak velká byla změna teploty? Vyjádřete tuto změnu pomocí absolutní hodnoty.

        \task \textbf{Měření chyby:}  
        Naměřená délka je $l_m = 24{,}8\,\mathrm{cm}$, skutečná délka $l_s = 25{,}0\,\mathrm{cm}$.  
        Určete absolutní chybu měření.

        \task \textbf{Vzdálenost na číselné ose:}  
        Jaká je vzdálenost bodů $A(-4)$ a $B(5)$?  
        Zapište výpočet pomocí absolutní hodnoty.

        \task \textbf{Pohyb po silnici:}  
        Auto jelo z bodu $x=0$ do bodu $x=-3$ a poté do bodu $x=5$.  
        Jakou celkovou dráhu ujelo? Zapište výpočet pomocí absolutní hodnoty.
    \end{tasks}
\end{exercise}


% =========================
% ÚLOHA 8
% =========================
\begin{exercise}[Graf funkce]
    Sestroj graf funkce: 
    \[
    m: y= 15 -2|x| - |2x-7| - 3\cdot |1+x|
    \]
    , pro $x\in \langle -3;5\rangle$. 
\end{exercise}


% =========================
% ÚLOHA 9
% =========================
\begin{exercise}
    Načrtněte graf funkce, pro kterou platí:
    \begin{tasks}
        \task Její definiční obor je $(-2;2)$, je lichá, je klesající.
        \task Její definiční obor je $\mathbb{R}$, je shora i zdola omezená, je sudá.
        \task Její definiční obor je $\langle0;5\rangle$, obor hodnot $\langle0;2\rangle$, v bodě $1$ má maximum, v bodech $0$ a $3$ je minimum.
    \end{tasks}
\end{exercise}



\end{document}